\documentclass[conference]{IEEEtran}
\IEEEoverridecommandlockouts

\usepackage{cite}
\usepackage{amsmath,amssymb}
\usepackage{graphicx}
\usepackage{booktabs}
\usepackage{url}
\usepackage[hidelinks]{hyperref}

\begin{document}

\title{Batch-Tiled Evaluation Overstates the Cost of\\Causal Sleep Staging: Evidence from Two Datasets}

\author{\IEEEauthorblockN{Jahnvi R, Swapnil S, Shreeya S Methuku, Shreshtha Dixit}
\IEEEauthorblockA{Department of Computer Science (AIML) and Engineering,
PES University, Bengaluru, India}}

\maketitle

\begin{abstract}
A sleep-staging model intended for a wearable device may use only past signal, and this
constraint is generally assumed to cost accuracy. We measure that cost with a matched ablation
in which one architecture is trained with and without access to future signal, holding
parameter count, data, folds and seeds fixed, and run the identical experiment on two datasets.
Under the batch-tiled protocol standard in this literature, which covers a recording with
non-overlapping context windows, the causal constraint costs $0.025$ Cohen's $\kappa$ on 78
Sleep-EDF subjects and $0.033$ on the 25 subjects of DOD-H, and on both appears concentrated in
REM. We show that most of this is produced by the protocol. Tiling scores predictions at every
position of the context buffer, including positions with almost no accumulated history, which a
streaming system passes through once at start-up; the causal model is crippled there while the
non-causal model is not, because it retains lookahead. Moving the same weights to a streaming
protocol therefore lifts the two arms unequally, and the excess lift to the causal arm is
$+0.0314$ and $+0.0292\,\kappa$, positive in all 30 paired folds across the two datasets, with
Nadeau--Bengio corrected $t(14) = 12.69$ and $6.48$. We verify the mechanism arithmetically: the
tiled score is the mean of a curve over buffer positions, reproduced to within $0.003\,\kappa$
in all four model--dataset combinations, and the penalty is generated by positions carrying
under 30 seconds of history. Re-measured under the deployment protocol the causality cost is
$+0.0065$ and $-0.0037\,\kappa$, neither resolvable, and the apparent REM concentration falls
from $0.060$ F1 to $0.002$ and $0.023$. Reported causality penalties in this literature should
be read together with the evaluation protocol that produced them.
\end{abstract}

\begin{IEEEkeywords}
sleep staging, EEG, causal models, real-time inference, evaluation methodology
\end{IEEEkeywords}

\section{Introduction}

Most automated sleep-staging models classify 30-second epochs and are free to use signal from
later in the recording, or the whole night, when labelling an earlier epoch~\cite{deepsleepnet,
attnsleep}. That is appropriate for offline analysis of a completed study but unsuitable for a
wearable device, which must report a stage as it occurs. Causality is therefore a real design
constraint, and the natural question is what it costs.

That question is rarely answered cleanly. Published causal and non-causal models differ in
capacity, data, preprocessing and optimization simultaneously, so the gap between them is not
attributable to causality in particular. High output resolution is not the open problem either:
U-Sleep~\cite{usleep} and its predecessor U-Time~\cite{utime} have emitted up to 7{,}680 stages
per minute since 2021, using a bidirectional model that reads several minutes of future signal.
What has not been measured is what that resolution costs when the model may not look forward at
all.

Causal staging itself is an active direction. Bresch et al.~\cite{bresch} apply recurrent
networks to real-time single-channel staging, CausalAttenNet~\cite{causalattennet} combines
causal convolution with long-term attention, and Micro SleepNet~\cite{microsleepnet} targets
wearable hardware. We claim no priority over these on the architecture. To our knowledge,
however, none reports a matched ablation in which the identical network is trained with and
without access to future signal on identical folds and seeds, which is what licenses attributing
a gap to causality rather than to capacity, and none reports how the answer moves with the
evaluation protocol.

We answer it with a matched ablation: one architecture, trained twice, differing only in whether
convolutions are left-padded and attention is masked, run end to end on two datasets. The first
result is the expected one. Under the evaluation protocol this literature normally uses,
causality costs $0.025\,\kappa$ on Sleep-EDF-78 and $0.033$ on DOD-H. The second result is that
most of that number is produced by the protocol rather than by the constraint, and this paper is
about the second result.

Our contribution is one finding with three consequences. The finding is that batch-tiled
evaluation systematically penalizes causal models, measured as a within-arm quantity that does
not require the two arms to be comparable, replicated on two datasets that differ in montage,
sampling rate, scoring team, class balance and subject count. The consequences are that the
reported causality penalty is largely an artifact, that its apparent concentration in REM is the
same artifact, and that a corpus too small to resolve the between-arm ablation can still resolve
this effect decisively, which is what makes the measurement worth adopting.

\section{Method}

\subsection{Datasets}

\textbf{Sleep-EDF-78}~\cite{physionet,kemp}: 78 subjects, 153 recordings, single channel Fpz-Cz
at 100~Hz, 1{,}629 hours after trimming each night to 30 minutes of Wake either side of the sleep
period. \textbf{DOD-H}~\cite{dod}: 25 subjects, 205.5 hours, channel F3-M2 resampled from 250 to
100~Hz with a polyphase anti-aliasing filter, no Wake trimming, since Wake is already 12.3\% of
the corpus against 68\% in untrimmed Sleep-EDF. Both are scored to five AASM classes~\cite{aasm}
in 30-second epochs; each epoch label is replicated across its 30 one-second targets, so the
network is supervised at 1~Hz. Both use subject-wise 5-fold cross-validation with no subject
crossing the train/test boundary. Nothing about the architecture, the training recipe or the two
protocols changes between datasets; the focal-loss class weights are recomputed per training
fold, so the differing class balances are handled automatically.

\subsection{Model and the Causality Ablation}

A multi-resolution causal CNN (kernels of 50 and 400 samples, left-aligned pooling to 1~Hz) feeds
three residual blocks of causal dilated convolutions~\cite{tcn} (dilations 1, 2, 4, 32 channels)
and a four-head self-attention layer~\cite{attention} with a causal mask, then a linear classifier
emitting one five-class prediction per second. Training uses focal loss~\cite{focal} with
$\gamma = 2$ and per-fold class weights, early stopping on validation $\kappa$. Total 30{,}757
parameters; 0.026~ms per second of EEG on one CPU thread.

Causality is enforced in exactly two places: convolutions pad only the left by $(K-1)d$, and the
attention mask sets $M_{ij} = -\infty$ for $j > i$. The non-causal arm removes both and changes
nothing else, so the two arms are identical in capacity and differ only in access to future
signal.

Causality also has to hold upstream of the network. Per-epoch $z$-scoring, near-universal in this
literature, normalizes an early sample using statistics over the whole epoch including its future.
We replace it with a trailing 30-second statistic computable one sample at a time. Closing this
leak raises $\kappa$ from $0.6406$ to $0.6649$ (causal) and $0.6693$ to $0.6898$ (non-causal), so
it is a modelling gain and never a trade. The full path from raw sample to logit is verified
causal by perturbation test, and a streaming front end reproduces the offline labels exactly over
a 79{,}440-second held-out recording. Because the two corpora are recorded in different amplitude
units, we also verify that the normalizer removes a $1000\times$ change in input gain to within
$5\times10^{-7}$.

\subsection{Evaluation Protocols}
\label{sec:protocol}

\textbf{Tiled.} The recording is covered by non-overlapping context windows and every position of
every window contributes one scored prediction. This is the default in every implementation we
have examined, including our own earlier work. Each window begins from an empty history.

\textbf{Streaming.} The model recomputes every 30~s against a full buffer and only the freshest 30
predictions are kept, so after the opening window every scored second carries at least 90 seconds
of genuine history. This is what a deployed system does. Every held-out second is scored exactly
once under each protocol, from the same trained weights.

\subsection{Statistics}

Each arm is trained under three seeds $\times$ five folds per dataset, giving 15 paired
measurements and 60 trained models in total. Because repeats reuse subjects, the primary test is
the Nadeau--Bengio corrected resampled $t$-test~\cite{nadeau}, which inflates the variance
estimate by $1/n + 1/(k-1) = 1/15 + 1/4$ here.
Bootstrap intervals resample whole subjects. Pairing is exact: within a seed the two arms use
byte-identical subject lists.

\section{Results}

\subsection{The Protocol Change Is Worth Ten Times More to the Causal Arm}

Table~\ref{tab:protocol} gives the causality cost under each protocol. On Sleep-EDF-78 the tiled
penalty is resolved and the streaming one is not. On DOD-H neither is: with five subjects per
held-out fold the corrected test returns $p = 0.22$ on a point estimate larger than
Sleep-EDF-78's. A claim built on this table alone would therefore rest on one dataset.

\begin{table}[t]
\caption{Cost of Causality Under Each Protocol. 15 Paired Folds Per Row, Nadeau--Bengio Corrected}
\label{tab:protocol}
\centering
\footnotesize
\setlength{\tabcolsep}{4pt}
\begin{tabular}{llccccc}
\toprule
Dataset & Protocol & Causal & Non-c. & $\Delta$ & Loses & $p$ \\
\midrule
Sleep-EDF-78 & Tiled     & 0.6649 & 0.6898 & $-0.0249$ & 14/15 & 0.005 \\
Sleep-EDF-78 & Streaming & \textbf{0.6998} & 0.6934 & $\mathbf{+0.0065}$ & 7/15 & 0.47 \\
\midrule
DOD-H        & Tiled     & 0.6329 & 0.6659 & $-0.0329$ & 10/15 & 0.22 \\
DOD-H        & Streaming & \textbf{0.6674} & 0.6711 & $\mathbf{-0.0037}$ & 6/15 & 0.89 \\
\bottomrule
\end{tabular}
\end{table}

The decisive statistic is a different one, and it is decisive because it does not require the two
arms to be comparable. Ask instead how much each arm gains from the protocol change, which is a
within-model difference on identical weights and identical held-out seconds
(Table~\ref{tab:asym}).

\begin{table}[t]
\caption{Protocol Gain by Arm: $\kappa$(Streaming) $-$ $\kappa$(Tiled), Paired Within Fold}
\label{tab:asym}
\centering
\footnotesize
\setlength{\tabcolsep}{5pt}
\begin{tabular}{lcc}
\toprule
 & Sleep-EDF-78 & DOD-H \\
\midrule
Causal arm gains       & $+0.0350$ & $+0.0345$ \\
Non-causal arm gains   & $+0.0036$ & $+0.0052$ \\
\midrule
\textbf{Excess to the causal arm} & $\mathbf{+0.0314}$ & $\mathbf{+0.0292}$ \\
Folds with positive excess & 15/15 & 15/15 \\
Bootstrap 95\% CI & $[.0291, .0334]$ & $[.0253, .0332]$ \\
Corrected $t(14)$ & $12.69$ & $6.48$ \\
$p$ & $4.6\times10^{-9}$ & $1.4\times10^{-5}$ \\
Per-seed range & $.0310$--$.0319$ & $.0285$--$.0297$ \\
\bottomrule
\end{tabular}
\end{table}

The protocol change is worth roughly ten times as much to the causal arm as to the non-causal
arm; the excess is positive in every one of the thirty paired folds across the two datasets, the
smallest single-fold value anywhere being $+0.0173$; and two corpora differing in electrode
montage, sampling rate, scoring team, class balance, subject count and recording length agree on
its size to within $0.002\,\kappa$. The aggregate cost follows by arithmetic rather than as an
independent observation: a tiled penalty of $-0.0249$ and $-0.0329$, an excess of $+0.0314$ and
$+0.0292$, and a streaming residual of $+0.0065$ and $-0.0037$. The protocol accounts for
essentially all of the measured penalty on Sleep-EDF-78 and about 89\% of it on DOD-H.

\subsection{The Mechanism, Verified}
\label{sec:latency}

Define permitted latency $L$: the label for second $t$ is emitted at $t+L$, so the model sees $L$
seconds of future and $L_{\text{buf}} - L$ of past, and $L = 0$ is real-time operation. One
forward pass over a window yields a prediction at every position, and position $p$ carries exactly
$L_{\text{buf}} - 1 - p$ seconds of lookahead, so a single sweep measures every latency.

Averaging the resulting curve over all 120 buffer positions reproduces that model's tiled score to
within $0.003\,\kappa$ in all four model--dataset combinations: $0.6626$ against a measured
$0.6634$ and $0.6890$ against $0.6912$ on Sleep-EDF-78, $0.6537$ against $0.6567$ and $0.6915$
against $0.6922$ on DOD-H. The tiled protocol is, arithmetically, this average.

Table~\ref{tab:quartile} locates the penalty, and the two datasets show the same shape. At buffer
positions carrying 91--120 seconds of history the arms are level. Across the last quartile the
causal arm loses $0.114$ and $0.128\,\kappa$ against its own best quartile while the non-causal arm
loses $0.018$ and $0.028$; that gap is the artifact, and a streaming system occupies those
positions once, at start-up.

\begin{table}[t]
\caption{$\kappa$ by Buffer Quartile, Seed 42, 120~s Buffer. Positive $\Delta$ Means the Causal Arm
Pays}
\label{tab:quartile}
\centering
\footnotesize
\setlength{\tabcolsep}{3.6pt}
\begin{tabular}{lccc|ccc}
\toprule
 & \multicolumn{3}{c|}{Sleep-EDF-78} & \multicolumn{3}{c}{DOD-H} \\
History (s) & Caus. & Non-c. & $\Delta$ & Caus. & Non-c. & $\Delta$ \\
\midrule
91--120 & \textbf{0.6982} & 0.6930 & $-0.0052$ & 0.6904 & 0.6961 & $+0.0057$ \\
61--90  & 0.6923 & 0.6969 & $+0.0046$ & 0.6887 & 0.7044 & $+0.0157$ \\
31--60  & 0.6753 & 0.6911 & $+0.0158$ & 0.6737 & 0.6970 & $+0.0233$ \\
1--30   & 0.5847 & 0.6752 & $+0.0905$ & 0.5622 & 0.6685 & $+0.1063$ \\
\bottomrule
\end{tabular}
\end{table}

Fig.~\ref{fig:latency} gives the full curves. The causal arm declines monotonically from $L = 0$:
it cannot read future signal, so larger $L$ only costs it history, and at $L = 119$ it collapses to
$\kappa = 0.317$ and $0.252$. The non-causal arm traces an inverted U and peaks at $L = 38$~s on
Sleep-EDF-78 and $L = 46$~s on DOD-H.

At each arm's best operating point the two datasets disagree, and we report that rather than the
agreement because it bounds the claim. On Sleep-EDF-78 the causal arm reaches $\kappa = 0.7008$ at
zero latency against a non-causal optimum of $0.6977$ at 38~s of delay, a margin of
$0.003\,\kappa$ we would not defend as more than a tie. On DOD-H the causal arm reaches $0.6922$ at
zero latency and the non-causal optimum is $0.7058$ at 46~s, exceeding it by $0.0136$. What replicates
across the two datasets is the protocol asymmetry; the stronger statement that a non-causal model
cannot beat a causal one at any latency holds on Sleep-EDF-78 and fails on DOD-H, where roughly
$0.014\,\kappa$ is available in exchange for 46 seconds of delay.

\begin{figure}[t]
\centering
\includegraphics[width=\columnwidth]{figures/fig_latency.png}
\caption{Cohen's $\kappa$ against permitted latency $L$, both arms, both datasets, 5-point moving
average over $L$. Larger $L$ means less history and more lookahead within a fixed 120~s buffer. The
causal collapse at large $L$ is the cold-start behaviour the tiled protocol averages into its
score.}
\label{fig:latency}
\end{figure}

\subsection{The REM Concentration Is the Same Artifact}

Under tiling, REM appears to pay $0.060$ F1 on both datasets, four times the next-largest class on
Sleep-EDF-78 and the only class whose penalty there survives Holm correction across the five
per-class tests. Under the deployment protocol the difference is $+0.002$ and $+0.023$, and N1, N2
and N3 invert on both datasets, the causal arm ending ahead on all three
(Table~\ref{tab:perclass}). This is mechanistically sensible: REM is defined partly by surrounding
context, which is why a technician resolves it against neighbouring epochs, so it is the class most
damaged by a nearly empty buffer. Two residuals remain rather than vanishing. REM still costs
$0.023$ F1 on DOD-H at deployment, and Wake costs $0.035$ there under both protocols; DOD-H is not
Wake-trimmed, so its Wake is sleep-adjacent and transitional throughout, which is exactly the case
where seeing forward helps.

\begin{table}[t]
\caption{Per-Class F1 Difference by Protocol, 15 Paired Folds Per Column. Positive Means the Causal
Arm Pays}
\label{tab:perclass}
\centering
\footnotesize
\setlength{\tabcolsep}{3.2pt}
\begin{tabular}{lcc|cc}
\toprule
 & \multicolumn{2}{c|}{Sleep-EDF-78} & \multicolumn{2}{c}{DOD-H} \\
Class & Tiled & Streaming & Tiled & Streaming \\
\midrule
W   & $+0.018$ & $+0.006$ & $+0.046$ & $+0.035$ \\
N1  & $+0.025$ & $-0.009$ & $+0.007$ & $-0.018$ \\
N2  & $+0.002$ & $-0.017$ & $+0.014$ & $-0.004$ \\
N3  & $+0.008$ & $-0.020$ & $+0.015$ & $-0.004$ \\
REM & $\mathbf{+0.060}$ & $\mathbf{+0.002}$ & $\mathbf{+0.060}$ & $\mathbf{+0.023}$ \\
\midrule
Macro & $+0.023$ & $-0.007$ & $+0.028$ & $+0.006$ \\
\bottomrule
\end{tabular}
\end{table}

\subsection{What the Per-Second Rate Buys}

A per-second rate is only useful if it lets a system report something sooner. Charging both systems
for their evidence, a change counts as reported once the new stage has been held for $h$ seconds,
while an epoch-level system cannot emit the label for epoch $k$ until epoch $k$ has ended. At
$h = 10$~s the per-second system reports first in 94.8\% of transitions on Sleep-EDF-78 and 95.6\%
on DOD-H, with a median 19 seconds saved on both. The advantage is conditional on $h$ and we report
the full sweep: it falls to 66.5\% at $h = 20$ and 29.6\% at $h = 30$, where the epoch system
leads, because requiring 30 seconds of confirmation reimposes the epoch and discards the
resolution.

\subsection{Why the Estimand Matters at This Sample Size}

DOD-H's fold-to-fold $\kappa$ standard deviation is $0.098$, four times Sleep-EDF-78's $0.026$,
because five subjects are held out per fold rather than sixteen. The same fifteen runs cannot
resolve the between-arm causality difference ($p = 0.22$) and do resolve the within-arm protocol
gain ($p = 1.4\times10^{-5}$), because the latter scores identical weights on identical seconds and
so removes model variance from the comparison entirely. At the subject counts common in this
literature the choice of estimand decides what is measurable, which is the practical reason to
report the protocol comparison alongside the ablation.

\section{Discussion and Conclusion}

An earlier version of this work reported that the causal constraint costs $0.029\,\kappa$ and that
the cost concentrates in REM. Both are reproducible under the protocol that produced them and both
are largely artifacts of it. The mechanism is not subtle once seen: tiling a recording with
independent context windows repeatedly destroys the accumulated history that defines a streaming
model and then scores it on the wreckage, and it does so asymmetrically, since a model that can
read forward is largely unaffected at the start of a window and a model that cannot is crippled
there.

Three things bound the claim. The streaming protocol denies the non-causal arm the lookahead that
defines it, so this compares deployment options rather than architectures; at matched latency the
answer is dataset-dependent, with lookahead worth nothing on Sleep-EDF-78 and about $0.014\,\kappa$
at 46 seconds of delay on DOD-H. The non-causal arm is our own matched twin rather than a
state-of-the-art offline model, so the comparison bounds what causality costs within this capacity
class. And both datasets are single-channel EEG scored to AASM conventions with one architecture;
the mechanism argues that any model whose predictions depend on accumulated history should show the
effect, but two instances of one architecture is not a general result.

What we would like adopted is smaller than a conclusion about causality and more robust than one.
Any matched ablation of a streaming model against a bidirectional one should report both protocols,
because on our data the two numbers differ by more than the effect being measured, and because the
within-arm comparison is resolvable on corpora where the between-arm one is not. Code, per-fold
predictions, splits, and the analysis for both datasets are released so the measurement can be
repeated without retraining.

\begin{thebibliography}{99}

\bibitem{deepsleepnet} A. Supratak, H. Dong, C. Wu, and Y. Guo, ``DeepSleepNet: A model for
automatic sleep stage scoring based on raw single-channel EEG,'' \emph{IEEE Trans. Neural Syst.
Rehabil. Eng.}, vol. 25, no. 11, pp. 1998--2008, 2017.

\bibitem{attnsleep} E. Eldele et al., ``An attention-based deep learning approach for sleep stage
classification with single-channel EEG,'' \emph{IEEE Trans. Neural Syst. Rehabil. Eng.}, vol. 29,
pp. 809--818, 2021.

\bibitem{usleep} M. Perslev, S. Darkner, L. Kempfner, M. Nikolic, P. J. Jennum, and C. Igel,
``U-Sleep: Resilient high-frequency sleep staging,'' \emph{npj Digit. Med.}, vol. 4, art. 72, 2021.

\bibitem{physionet} A. L. Goldberger et al., ``PhysioBank, PhysioToolkit, and PhysioNet: Components
of a new research resource for complex physiologic signals,'' \emph{Circulation}, vol. 101, no. 23,
pp. e215--e220, 2000.

\bibitem{kemp} B. Kemp, A. H. Zwinderman, B. Tuk, H. A. C. Kamphuisen, and J. J. L. Oberye,
``Analysis of a sleep-dependent neuronal feedback loop: the slow-wave microcontinuity of the EEG,''
\emph{IEEE Trans. Biomed. Eng.}, vol. 47, no. 9, pp. 1185--1194, 2000.

\bibitem{dod} A. Guillot, F. Sauvet, E. H. During, and V. Thorey, ``Dreem Open Datasets:
Multi-scored sleep datasets to compare human and automated sleep staging,'' \emph{IEEE Trans.
Neural Syst. Rehabil. Eng.}, vol. 28, no. 9, pp. 1955--1965, 2020.

\bibitem{aasm} R. B. Berry et al., \emph{The AASM Manual for the Scoring of Sleep and Associated
Events, Version 2.6}. Darien, IL: American Academy of Sleep Medicine, 2020.

\bibitem{tcn} S. Bai, J. Z. Kolter, and V. Koltun, ``An empirical evaluation of generic
convolutional and recurrent networks for sequence modeling,'' \emph{arXiv:1803.01271}, 2018.

\bibitem{attention} A. Vaswani et al., ``Attention is all you need,'' in \emph{Advances in Neural
Information Processing Systems 30}, 2017, pp. 5998--6008.

\bibitem{focal} T.-Y. Lin, P. Goyal, R. Girshick, K. He, and P. Doll\'ar, ``Focal loss for dense
object detection,'' in \emph{Proc. IEEE Int. Conf. Comput. Vis. (ICCV)}, 2017, pp. 2980--2988.

\bibitem{nadeau} C. Nadeau and Y. Bengio, ``Inference for the generalization error,'' \emph{Mach.
Learn.}, vol. 52, no. 3, pp. 239--281, 2003.

\bibitem{causalattennet} J. Pan, Y. Feng, P. Zhao, X. Zou, A. Hou, and X. Che, ``CausalAttenNet: A
fast and long-term-temporal network for automatic sleep staging with single-channel EEG,''
\emph{IEEE Trans. Instrum. Meas.}, vol. 73, pp. 1--13, 2024.

\bibitem{bresch} E. Bresch, U. Gro\ss{}ekath\"ofer, and G. Garcia-Molina, ``Recurrent deep neural
networks for real-time sleep stage classification from single channel EEG,'' \emph{Front. Comput.
Neurosci.}, vol. 12, art. 85, 2018.

\bibitem{microsleepnet} G. Liu et al., ``Micro SleepNet: Efficient deep learning model for mobile
terminal real-time sleep staging,'' \emph{Front. Neurosci.}, vol. 17, art. 1218072, 2023.

\bibitem{utime} M. Perslev, M. Jensen, S. Darkner, P. J. Jennum, and C. Igel, ``U-Time: A fully
convolutional network for time series segmentation applied to sleep staging,'' in \emph{Advances in
Neural Information Processing Systems 32}, 2019.

\end{thebibliography}

\end{document}
